@ID: ON22D1P1U
@BIDANG: Struktur Aljabar
Misalkan $a, b,$ dan $c$ adalah unsur di grup $G$ sehingga $abc=e$ dengan $e$ adalah unsur identitas di $G$. Untuk masing-masing pernyataan berikut, buktikan jika benar atau berikan contoh penyangkal jika salah.
(a) $bca=e$
(b) $bac=e$

@ID: ON22D1P2U
@BIDANG: Kombinatorika
Misalkan $n\in\mathbb{N}$ dan $c$ adalah sembarang bilangan real. Tunjukkan bahwa terdapat bilangan bulat $a$ dan $b$ dengan $1\le b\le n$ sedemikian sehingga $|c-\frac{a}{b}|\le\frac{1}{ab}$.

@ID: ON22D1P3U
@BIDANG: Analisis Kompleks
Diberikan suku banyak $p(z)=az^{2}+bz+c$ untuk suatu bilangan kompleks $a, b,$ dan $c$.
(a) Jika $\{\omega_{0},\omega_{1},\omega_{2}\}$ merupakan himpunan penyelesaian dari persamaan $\omega^{3}=1$, tentukan nilai dari

$$
\omega_{0}p(\omega_{0})+\omega_{1}p(\omega_{1})+\omega_{2}p(\omega_{2})
$$

(b) Jika $|abc|>1$, tunjukkan bahwa terdapat bilangan kompleks $z$ dengan $|z|\le1$ sehingga $|p(z)|>1$.

@ID: ON22D1P4U
@BIDANG: Aljabar Linear
Diberikan matriks $A$ berukuran $4\times2$ dan matriks

$$
B=\begin{bmatrix}b_{11}&b_{12}&b_{13}&b_{14}\\ b_{21}&b_{22}&b_{23}&b_{24}\end{bmatrix}
$$

yang memenuhi

$$
A\begin{bmatrix}b_{11}\\ b_{21}\end{bmatrix}=\begin{bmatrix}0\\ 2\\ 0\\ 0\end{bmatrix}
$$

dan

$$
A\begin{bmatrix}b_{12}+b_{13}+b_{14}\\ b_{22}+b_{23}+b_{24}\end{bmatrix}=\begin{bmatrix}2\\ 0\\ 2\\ 2\end{bmatrix}
$$

(a) Tentukan semua nilai eigen dari matriks $AB$ beserta multiplisitas aljabarnya.
(b) Buktikan bahwa

$$
BA=\begin{bmatrix}2&0\\ 0&2\end{bmatrix}
$$

@ID: ON22D1P5U
@BIDANG: Analisis Real
Misalkan barisan bilangan real $(a_{n})$ dengan $\frac{1}{2}<a_{n}<1$ untuk setiap $n\ge0$. Didefinisikan barisan $(x_{n})$ dengan

$$
x_{0}=0, x_{n+1}=\frac{a_{n+1}+x_{n}}{1+a_{n+1}x_{n}}, n\ge0
$$

Apakah $(x_{n})$ konvergen? Jika ya, tentukan nilai limitnya.

@ID: ON22D2P1U
@BIDANG: Kombinatorika
Misalkan $k\ge4$ adalah bilangan bulat genap dan $S=\{1,2,...,n\}$. Tentukan bilangan asli terkecil $n$ sedemikian sehingga untuk setiap pewarnaan bilangan-bilangan di dalam $S$ dengan tiga warna berbeda, senantiasa terdapat tiga anggota $S$, katakan $a, b,$ dan $c$ (tidak harus berbeda) yang berwarna sama dan memenuhi $ka+b=c$.

@ID: ON22D2P2U
@BIDANG: Analisis Real
Diberikan fungsi terbatas $f:[a,b]\rightarrow\mathbb{R}$ dan $x_{0}\in(a,b)$. Didefinisikan fungsi $p:(a,b)\rightarrow\mathbb{R}$ dengan

$$
p(x)=\sup\{f(t):t<x\}
$$

Jika fungsi $f$ kontinu di $x_{0}$, selidiki apakah $p$ kontinu di $x_{0}$. Berikan penjelasan jawaban saudara.

@ID: ON22D2P3U
@BIDANG: Analisis Kompleks
Diberikan fungsi kompleks

$$
f(z)=1+c_{1}z+c_{2}z^{2}+\dots+c_{k}z^{k}
$$

dengan $|c_{n}|<1945$ untuk setiap $1\le n\le k$. Buktikan bahwa fungsi $f$ tidak memiliki akar-akar $z$ dengan sifat

$$
|z|<\frac{1}{2022}
$$

@ID: ON22D2P4U
@BIDANG: Struktur Aljabar
Misalkan $A$ adalah ring dengan identitas perkalian $1$. Untuk setiap $a\in A$, didefinisikan pemetaan $l_{a}$ dan $r_{a}$ dari $A$ ke $A$ melalui $l_{a}(x)=ax$ dan $r_{a}(x)=xa$ untuk setiap $x\in A$.

(a) Berikan contoh ring tak komutatif $A$ sehingga manakala $a\in A$ membuat salah satu dari $l_{a}$ atau $r_{a}$ injektif, maka yang satunya lagi juga injektif.

(b) Berikan contoh ring tak komutatif $A$ sehingga terdapat $a\in A$ yang membuat tepat salah satu diantara $l_{a}$ atau $r_{a}$ injektif.

@ID: ON22D2P5U
@BIDANG: Aljabar Linear
Diberikan dua bilangan bulat positif $n$ dan $m$ dengan $n\le m$. Buktikan bahwa

$$
\begin{vmatrix}1&_{m}P_{1}&\dots&_{m}P_{n}\\ 1&_{m+1}P_{1}&\dots&_{m+1}P_{n}\\ \vdots&\vdots&\ddots&\vdots\\ 1&_{m+n}P_{1}&\dots&_{m+n}P_{n}\end{vmatrix}=1!2!\dots n!
$$

Catatan: $|A|$ menyatakan determinan dari matriks $A$ dan

$$
_{r}P_{s}:=r(r-1)(r-2)\dots(r-s+1)
$$

untuk setiap $r, s\in\mathbb{N}$ dan $s\le r$.